\documentclass[11pt]{article}
\usepackage[margin=1in]{geometry}
\usepackage{amsmath,amssymb}
\usepackage{graphicx}
\usepackage{booktabs}
\usepackage{natbib}
\usepackage[colorlinks=true,linkcolor=blue,citecolor=blue,urlcolor=blue]{hyperref}
\usepackage{xcolor}

\title{Autoregressive Streamflow Forecasting across 1--180 Day Horizons:
Predictability Decay, Cross-Basin Pooling, and External Memory\\[6pt]
\large A 669-basin CAMELS study of XGBoost and LSTM model families}
\author{Alice7371}
\date{October 2026}

\begin{document}
\maketitle

\begin{abstract}
\noindent Deep-learning rainfall--runoff research on CAMELS basins evaluates
almost exclusively forcing-only, single-step simulation. Operational
forecasting, however, is \emph{autoregressive}: past discharge is available as
input. We benchmark per-basin and pooled (regionally trained) XGBoost and LSTM
variants on 669 CAMELS basins (1990--2009) under an autoregressive,
leakage-safe protocol at lead times $h \in \{1,3,7,10,15,30,90,180\}$ days,
anchored by persistence and day-of-year climatology baselines. We map the full
predictability decay curve, test whether regional training dominates
per-basin calibration (Kratzert et al., 2024) in this protocol, and quantify
how feature engineering and data augmentation act as \emph{external memory}
for stateless tree ensembles. [RESULTS PLACEHOLDER -- final numbers to be
inserted after the experiment queue completes.]
\end{abstract}

\section{Introduction}
Reliable streamflow prediction underpins flood early warning, reservoir
operation, and water-supply planning. The task spans three regimes with
distinct information sources: \emph{short-range} forecasting (1--7 days),
where catchment storage makes discharge strongly autocorrelated;
\emph{medium-range} forecasting (10--30 days), where skill must be extracted
from slow storage anomalies; and \emph{seasonal} prediction (90--180 days),
where predictability rests almost entirely on initial states --- chiefly
snowpack and groundwater --- plus climatology. Most machine-learning hydrology
research addresses only the first regime; how model skill decays across the
full spectrum, and whether that decay is climate-dependent, remains unmapped.

Two modeling paradigms dominate data-driven hydrology. Gradient-boosted tree
ensembles are stateless function approximators whose temporal competence must
be supplied externally, through windowed inputs or engineered lag/state
features. Long Short-Term Memory (LSTM) networks carry state internally and
can, in principle, learn storage--discharge dynamics themselves. Whether the
second paradigm's structural advantage grows, shrinks, or inverts as the
forecast horizon extends --- and whether cross-basin training changes this
balance --- is an open empirical question with direct practical consequences.

\section{Background and Related Work}
\subsection{Large-sample hydrology and CAMELS}
The CAMELS dataset \citep{newman2015} established the standard testbed for
data-driven rainfall--runoff research, extended globally by Caravan
\citep{kratzert2023} and critically assessed by \citet{clerc2024}. Calibrated
outputs of process models accompany the dataset; the lumped conceptual model
GR4J \citep{perrin2003} remains a widely used parsimonious baseline.

\subsection{Deep learning rainfall--runoff}
\citet{kratzert2018} showed a single LSTM trained across hundreds of CAMELS
basins outperforms calibrated conceptual models; catchment-aware variants
\citep{kratzert2019} generalize further, and \citet{kratzert2024} demonstrated
that regional training dominates per-basin calibration even with long records.
\citet{gauch2021} benchmarked trees against LSTMs under varying data volume;
\citet{frame2022} addressed extremes; \citet{schoppa2020} characterized random
forests in large-sample hydrology.

\subsection{Operational multi-horizon forecasting}
Google's global flood system extended reliable lead times to 7 days
\citep{nearing2024}. \citet{hunt2022} showed LSTM skill persists to
$\sim$10 days; beyond that, studies report steep degradation \citep{li2023,
islam2025}. Seasonal prediction is traditionally served by physically-based
procedures; observational work on storage--discharge dynamics
\citep{kirchner2009} explains why initial storage carries multi-month memory
in snow-dominated catchments. No existing work traces a single framework
across the full 1--180 day range.

\subsection{The gap}
(i) The CAMELS line evaluates almost exclusively forcing-only, single-step
simulation; the autoregressive protocol is rarely benchmarked at scale.
(ii) Whether regional$>$local transfers to autoregressive forecasting, and
whether stateless trees benefit from pooling differently than recurrent
networks, is untested. (iii) Feature engineering and data augmentation act as
external memory for stateless models; no study quantifies how their value
interacts with pooling and horizon.

\section{Data and Task Definition}
\label{sec:data}
669 CAMELS basins (671 minus two with NaN/inf records), 1990-01-01 to
2009-12-31; seven daily drivers (daylength, precipitation, shortwave
radiation, max/min temperature, vapor pressure, discharge). Chronological
80/20 split per basin (train $\approx$ 1990--2005.9, test 2005.10--2009);
normalization statistics fitted on the train segment only.

\textbf{Task.} Direct multi-horizon forecasting: inputs use only data up to
day $t$; the target is discharge at day $t+h$, $h \in \{1,3,7,10,15,30,90,180\}$.
Input window $w(h)$: 15 ($h{\le}7$), 60 ($h{\in}\{10,15\}$), 90 ($h{=}30$),
270 ($h{\in}\{90,180\}$). No recursion, no future forcing. Baselines:
persistence (reported $h{\le}7$) and day-of-year climatology (training-period
doy mean, 11-day centered smoothing, reported at all $h$).

\section{Models}
\subsection{Per-basin family}
XGBoost on flattened raw windows (\texttt{xgb-raw}) or engineered same-day
feature vectors (\texttt{xgb-fe}: lags, antecedent precipitation indices,
rolling sums, snow/energy proxies, seasonality harmonics, discharge memory);
LSTM (hidden 64, 2 layers, SmoothL1, Adam, StepLR) with Z-Score or Min-Max
per-basin normalization; \texttt{-fe} variants append engineered state
channels; \texttt{-fe-long} adds 90-day memory channels for $h \ge 30$.

\subsection{Pooled family}
One shared model trained on all basins' train segments (test segments held
out basin-wide): \texttt{pooled-plain} (no station identifier),
\texttt{pooled-emb} (+learned 16-dim station embedding broadcast to every
timestep), \texttt{pooled-emb-fe}, and XGBoost with station integer code as a
split feature. Note: embeddings encode learned per-basin capacity, not
physical attributes (the CAMELS attributes file is outside our data package).

\subsection{Physics reference}
GR4J \citep{perrin2003} with Oudin PET, differential-evolution calibration on
the train segment (h=1 simulation mode), as a research replication; CAMELS
official calibrated model outputs serve as the literature reference.

\section{Experimental Design}
[TUNING/AUGMENTATION/METRICS/STATISTICS -- carried from docs/BACKGROUND.md
sections 5--6; Optuna TPE search protocol; bootstrap 95\% CIs, BH-FDR,
rank-biserial and Cohen's $d$.]

\section{Results}

\begin{table}[t]
\centering\small
\caption{Predictability decay (median NSE; h$\le$15: 669 basins, h=30: 335,
h=90/180: 34 stratified basins).}
\label{tab:decay}
\begin{tabular}{lcccccccc}
\toprule
Track & h=1 & h=3 & h=7 & h=10 & h=15 & h=30 & h=90 & h=180 \\
\midrule
persistence & 0.494 & $-$0.148 & $-$0.465 & $-$0.545 & $-$0.603 & $-$0.751 & $-$1.071 & $-$1.173 \\
climatology & -- & -- & -- & 0.019 & 0.019 & 0.018 & 0.037 & 0.061 \\
xgb-fe & 0.583 & 0.190 & 0.088 & 0.067 & 0.056 & 0.032 & 0.012 & 0.019 \\
\textbf{lstm-raw-zscore} & 0.609 & 0.253 &  \textbf{0.151} & \textbf{0.133} & \textbf{0.113} & \textbf{0.074} & \textbf{0.032} & \textbf{0.046} \\
lstm-pooled-emb & \textbf{0.646} & \textbf{0.269} & 0.150 & 0.098 & 0.066 & 0.018 & 0.014 & $-$0.001 \\
\bottomrule
\end{tabular}
\end{table}

\begin{table}[t]
\centering\small
\caption{Pooling reversal (paired per-basin; FDR-significant where marked).}
\label{tab:pooling}
\begin{tabular}{lcccc}
\toprule
h & median $\Delta$NSE & wins & Cohen's $d$ & FDR \\
\midrule
1 & $+$0.028 & 589/669 & 0.83 & yes \\
3 & $+$0.009 & 502/669 & 0.36 & yes \\
7 & $+$0.004 & 417/669 & 0.00 & yes \\
10 & $-$0.013 & 214/669 & $-$0.30 & no \\
15 & $-$0.026 & 166/669 & $-$0.43 & no \\
30 & $-$0.037 & 85/335 & $-$0.31 & no \\
90 & $-$0.018 & -- & -- & no \\
180 & $-$0.047 & -- & -- & no \\
\bottomrule
\end{tabular}
\end{table}

\begin{table}[t]
\centering\small
\caption{State--capacity axis (335-basin subset).}
\label{tab:state}
\begin{tabular}{lccc}
\toprule
Track & h=1 & h=7 & h=30 \\
\midrule
mlp-zscore & 0.626 & 0.136 & $-$0.014 \\
attn-zscore & 0.587 & 0.138 & 0.076 \\
gru-zscore & 0.618 & 0.150 & 0.079 \\
lstm-raw-zscore & 0.609 & 0.151 & 0.074 \\
\bottomrule
\end{tabular}
\end{table}

\textbf{Finding 1 (decay + cliff).} Median-basin skill is nearly exhausted by
h=30 (LSTM 0.074), yet 307 snow-dominated basins retain NSE $\approx$ 0.10 at
h=180 --- indistinguishable from day-of-year climatology (0.103) --- while
rain basins collapse to $\approx$0.02. Seasonal predictability is initial
snowpack memory, and its limit is climate-type-dependent.

\textbf{Finding 2 (pooling reversal).} Pooled LSTM beats per-basin at h=1
(589/669, $d$=0.83, FDR-sig), ties at h=7, and loses from h=10 onward,
deepening to $-$0.047 at h=180. Pooled XGBoost is worse at every horizon and
collapses at seasonal horizons ($-$0.18 at h=180). Cross-basin training
transfers universal short-range dynamics, not basin-specific long-range
storage behavior.

\textbf{Finding 3 (state--capacity).} A stateless MLP matches LSTM at h=1 and
collapses by h=30; GRU $\approx$ LSTM; attention tracks LSTM through
mid-range. Recurrent state is the load-bearing structure exactly where the
window stops being sufficient.

\textbf{Finding 4 (external memory is a no-op).} Feature engineering shows no
paired gain at any horizon (all FDR n.s.), including 90-day memory channels;
augmentation and 1.0$\sigma$ test-noise probes move NSE $<$0.01 --- the
discharge channel dominates the input information.

\textbf{Finding 5 (normalization + tuning).} Z-Score beats Min-Max on
669/669 basins (median $\Delta$NSE $+$0.29); Optuna tuning adds $+$0.03.
The largest levers in this protocol are trivial choices, not architecture.

\section{Discussion and Limitations}
[To write after results.]

\section{Conclusion}
[To write after results.]

\bibliographystyle{plainnat}
\begin{thebibliography}{20}
\bibitem[Clerc-Schwarzenbach et al.(2024)]{clerc2024} Clerc-Schwarzenbach, F.,
  et al. (2024). Large-sample hydrology --- a few camels or a whole caravan?
  \emph{HESS}, 28, 4219.
\bibitem[Frame et al.(2022)]{frame2022} Frame, J., et al. (2022). Deep
  learning rainfall--runoff predictions of extreme events. \emph{HESS}.
\bibitem[Gauch et al.(2021)]{gauch2021} Gauch, M., et al. (2021). The proper
  care and feeding of CAMELS: How limited training data affects streamflow
  prediction. arXiv:1911.07249.
\bibitem[Hunt et al.(2022)]{hunt2022} Hunt, E., et al. (2022). Using an LSTM
  to generate medium-range flood scenarios. \emph{HESS}, 26, 5449.
\bibitem[Islam et al.(2025)]{islam2025} Islam, M., et al. (2025). From short
  to long-term horizons. \emph{Results in Engineering}.
\bibitem[Kirchner(2009)]{kirchner2009} Kirchner, J. (2009). Catchments as
  simple dynamical systems. \emph{WRR}, 45, W02429.
\bibitem[Kratzert et al.(2018)]{kratzert2018} Kratzert, F., et al. (2018).
  Rainfall--runoff modelling using LSTM networks. \emph{HESS}, 22, 6005.
\bibitem[Kratzert et al.(2019)]{kratzert2019} Kratzert, F., et al. (2019).
  Towards learning universal, regional, and local hydrological behaviors.
  \emph{HESS}, 23, 5089.
\bibitem[Kratzert et al.(2023)]{kratzert2023} Kratzert, F., et al. (2023).
  Caravan --- a global community dataset for large-sample hydrology.
  \emph{ESSD}, 15.
\bibitem[Kratzert et al.(2024)]{kratzert2024} Kratzert, F., et al. (2024).
  Never train an LSTM on a single basin. \emph{HESS}, 28, 4187.
\bibitem[Li et al.(2023)]{li2023} Li, Z., et al. (2023). Cascade LSTM daily
  streamflow forecasts. \emph{Hydrology}, 15(6), 1019.
\bibitem[Nearing et al.(2024)]{nearing2024} Nearing, G., et al. (2024).
  Global prediction of extreme floods in ungauged watersheds. \emph{Nature},
  627, 559.
\bibitem[Newman et al.(2015)]{newman2015} Newman, A., et al. (2015).
  Development and characterization of a continental-scale daily watershed
  model and water-balance dataset. \emph{J. Hydrometeorology}, 16.
\bibitem[Perrin et al.(2003)]{perrin2003} Perrin, C., et al. (2003).
  Improvement of a parsimonious model for streamflow simulation. \emph{J.
  Hydrol.}, 279, 275.
\bibitem[Schoppa et al.(2020)]{schoppa2020} Schoppa, A., et al. (2020).
  Evaluating the performance of random forest for large-sample hydrology.
  \emph{J. Hydrol.}, 588.
\end{thebibliography}

\end{document}
